\documentclass[a4paper]{article}
\usepackage[pdftex]{graphicx}
\usepackage[utf8]{inputenc}
\usepackage{enumerate}
\usepackage{siunitx}
\sisetup{locale=DE} 
\usepackage{href-ul}
\hypersetup{
	colorlinks=true,
	linkcolor=blue,
	urlcolor=blue}
\usepackage{icomma}
\usepackage{amssymb}
\usepackage{geometry}
\geometry{a4paper, top=15mm, left=15mm, right=15mm, bottom=15mm,
	headsep=10mm, footskip=12mm}

% Start the document
\begin{document}
%Titel
{\bf Stegreifaufgabe aus der Physik }\\

\noindent Ein Kügelchen der Masse $m=\SI{12}{\milli\gram}$ und der Ladung $Q=\SI{-1,4e-9}{\coulomb}$ befindet sich in Ruhe auf einer horizontalen Unterlage und wird zum Zeitpunkt $t_0=\SI{0}{\second}$ in der Mitte des homogenen Feldes eines Plattenkondensators losgelassen. Die elektrische Feldstärke beträgt 
$\SI{36}{\kilo\volt\over \meter}$. Reibungsphänomene können vernachlässigt werden.

\begin{enumerate}[1.]

\item Erklären Sie, warum beim Laden des Kondensators kein konstanter Ladestrom fließt.

\item Zeigen Sie durch allgemeine Rechnung, dass für die Geschwindigkeit, mit welcher das Kügelchen auf die Kondensatorplatte prallt, in Abhängigkeit vom Plattenabstand $d$ gilt:
$$ v(d)=\SI{2,0}{\sqrt{\meter}\over \second}\cdot\sqrt{d}$$

\item Berechnen Sie diejenige Spannung $U$, die am Kondensator angelegt werden muss, damit das Kügelchen mit der Geschwindigkeit $v=\SI{0,56}{\meter\over \second}$ auf die Kondensatorplatte prallt.

\item Verändert sich diese Aufprallgeschwindigkeit, falls der Plattenabstand bei bestehender Verbindung zur Spannungsquelle vergrößert wird? Begründen Sie Ihre Antwort.

\item Was geschieht mit dem Kügelchen, nachdem es auf die Kondensatorplatte aufgeprallt ist? 

\end{enumerate} 

\fbox{
	\begin{minipage}{0.5\textwidth}
		Zur Lösung bitte \href{https://www.okuyakl.de/phys/p12efeldL093/ll093.pdf}{hier klicken} oder den QR-Code scannen.\\
	Weitere Arbeitsblätter gibt es unter 
	
	\href{https://www.okuyakl.de}{www.okuyakl.de}
	\end{minipage}
	\hfill
	\begin{minipage}{0.4\textwidth}
		\includegraphics[width=1.5 cm]{../../viecher/zwe03}
		\includegraphics[width=3 cm]{qr093}
		\includegraphics[width=2 cm]{../../viecher/afanticon1}
		
	\end{minipage}}
\end{document}
